\FloatBarrier
\section{The Block Power Method}
\label{sec:numerics}
\label{sec:nonhermitian}

Section~\ref{sec:transverse} reduced the dynamics to the leading eigenpairs of the transfer matrix $\mathcal{E}$. This section presents the algorithm we use to obtain them: a block power method that iterates several temporal states at once and returns the dominant eigenvalues together with their left and right eigenvectors, identified and ordered for the comparison with the \ac{CFT} predictions. The implementation details are collected in Appendix~\ref{app:blockpm}.

The right and left eigenvectors of $\mathcal{E}$ satisfy
\begin{equation}
  \mathcal{E}\ket{R_i}=\mu_i\ket{R_i},
  \qquad
  \bra{L_i}\mathcal{E}=\mu_i\bra{L_i}.
  \label{eq:nonhermitian_eigenproblem}
\end{equation}
Because $\mathcal{E}$ is non-normal (Section~\ref{sec:mporep:fsm}), the left and right eigenvectors are different objects, and neither set is orthogonal on its own; instead, the two sets are biorthogonal, $\braket{L_i}{R_j}\propto\delta_{ij}$. The bra notation involves no conjugation: $\bra{L_i}$ is a row vector, and $\braket{L_i}{R_j}$ is the direct contraction of the two temporal networks. Numerically, the left states are propagated with the transpose $\mathcal{E}^{\mathsf T}$, which is the matrix form of the left equation above.

The dominant eigenvalue $\mu_0$ carries the Loschmidt rate in its modulus and the central charge in its phase, the gaps to the subleading eigenvalues carry the boundary operator content, and the generalized temporal entropies are built from the dominant pair $(\bra{L_0},\ket{R_0})$. The calculation must therefore resolve the leading eigenvalues individually, identify $\mu_0$ among them, and order the rest so that they can be matched to the conformal tower, as we show below.

\subsection{The algorithm}
\label{sec:blockpm}

The dominant eigenpair alone is within reach of a single-vector power method, and for the entropy route, the Loschmidt rate and the central charge from the phase of $\mu_0$, that cheaper pipeline suffices. The boundary operator content, however, lives in the subleading eigenvalues, and their separations close as $T$ grows. To resolve them individually we iterate a block of $k$ states, which converges to the leading $k$-dimensional invariant subspace at a rate set by $|\mu_{k+1}/\mu_k|$, so near-degenerate eigenvalues inside the block do not limit it. The members are separated afterwards by the projected eigenproblem below.

The block method evolves $k$ right \acp{tMPS}, $\{\ket{R_1},\dots,\ket{R_k}\}$, together with $k$ left \acp{tMPS}, $\{\bra{L_1},\dots,\bra{L_k}\}$. The two blocks are initialized independently with random complex tensors. Algorithm~\ref{alg:blockpm} summarises one iteration, while Figure~\ref{fig:blockpm} shows the tensor-network operations involved in applying and projecting the transfer matrix.

\begin{figure}[htbp]
  \centering
  \includegraphics[width=0.8\textwidth]{imgs/blockpm.pdf}
  \caption{\textbf{One iteration of the block power method.} Each \ac{tMPS} is a column, and $\mathcal{E}$ is the corresponding column of the space--time network, closed at both ends by the product state (small open circles). (a) Application of $\mathcal{E}$ to one member of the right block, Eq.~\eqref{eq:blockpm_apply}: the temporal bond dimension grows from $\chi$ to $\chi \cdot D$, with $D$ the \ac{tMPO} bond dimension, and is truncated back to $\chi$. The left block is treated analogously with the transpose. (b) The contractions defining the two $k\times k$ matrices of Eq.~\eqref{eq:blockpm_pencil}; the left chain enters directly as a bra, without conjugation.}
  \label{fig:blockpm}
\end{figure}

\begin{algorithm}[H]
\caption{\textbf{Block power method} for the $k$ dominant eigenpairs of $\mathcal{E}$.}
\label{alg:blockpm}
\begin{algorithmic}[1]
\State $\ket{R_j},\ \bra{L_j} \gets$ random complex \acp{tMPS}, \quad $j=1,\dots,k$
\Repeat
  \State $\ket{\tilde R_j} \gets \operatorname{truncate}_\chi\big(\mathcal{E}\ket{R_j}\big)$, \qquad $\bra{\tilde L_j} \gets \operatorname{truncate}_\chi\big(\bra{L_j}\,\mathcal{E}\big)$ \Comment{apply}
  \State $S_{ij} \gets \braket{L_i}{R_j}$, \qquad $M_{ij} \gets \braket{L_i}{\tilde R_j}=\bra{L_i}\mathcal{E}\ket{R_j}$ \Comment{project}
  \State solve $M v=\theta\,S v$ for Ritz values $\theta_1,\dots,\theta_k$ and coefficients $V$, $U$
  \State $\ket{R_j} \gets \operatorname{truncate}_\chi\big(\textstyle\sum_i V_{ij}\ket{\tilde R_i}\big)$, \qquad $\bra{L_j} \gets \operatorname{truncate}_\chi\big(\textstyle\sum_i U_{ij}\bra{\tilde L_i}\big)$ \Comment{de-mix}
  \State normalize every block member
\Until{the leading Ritz values stop changing within tolerance}
\end{algorithmic}
\end{algorithm}

The first step applies the transfer matrix to each member of the right block and its transpose to each member of the left block,
\begin{equation}
  \ket{\tilde R_j} = \mathcal{E}\ket{R_j},
  \qquad
  \bra{\tilde L_j} = \bra{L_j}\,\mathcal{E}
  \;\;\Longleftrightarrow\;\;
  \ket{\tilde L_j} = \mathcal{E}^{\mathsf T}\ket{L_j}.
  \label{eq:blockpm_apply}
\end{equation}
Applying the \ac{tMPO} multiplies the bond dimension of each state by the \ac{tMPO} bond dimension $D$, as shown in Figure~\ref{fig:blockpm}a. We therefore truncate each propagated state back to the maximum bond dimension $\chi$ using either scheme of Section~\ref{sec:transverse:truncation}, and then normalize it; truncating after every application is what keeps the cost of the iteration bounded.

The next step projects the eigenproblem onto the current block subspaces. We approximate a right eigenvector by
\begin{equation}
  \ket{\psi} = \sum_{j=1}^{k} v_j \ket{R_j}.
  \label{eq:blockpm_ansatz}
\end{equation}
This approximation leaves a residual when inserted into $\mathcal{E}\ket{\psi}=\theta\ket{\psi}$. Requiring the residual to vanish against every member of the left block gives $\bra{L_i}(\mathcal{E}-\theta)\ket{\psi}=0$. The full eigenproblem is thereby reduced to the two $k\times k$ matrices shown in Figure~\ref{fig:blockpm}b,
\begin{equation}
  S_{ij}=\braket{L_i}{R_j},
  \qquad
  M_{ij}=\bra{L_i}\mathcal{E}\ket{R_j},
  \label{eq:blockpm_pencil}
\end{equation}
and the generalized eigenvalue problem
\begin{equation}
  \sum_{j} M_{ij}\,v_j = \theta \sum_{j} S_{ij}\,v_j .
  \label{eq:blockpm_smallproblem}
\end{equation}
The overlap matrix $S$ measures the relative orientation of the left and right subspaces, while $M$ is the transfer matrix projected between them. Solving Eq.~\eqref{eq:blockpm_smallproblem} returns the Ritz values $\theta_j$ together with two sets of coefficients. The right eigenvectors $V$ of the $k\times k$ pencil say how to combine the propagated right states into one approximate eigenvector each, and the left eigenvectors $U$ do the same for the left block. These are the $V$ and $U$ of Algorithm~\ref{alg:blockpm}. The resulting Ritz values approximate $\mu_0,\dots,\mu_{k-1}$, and the coefficient matrices identify the corresponding combinations within each block. If $S$ becomes ill-conditioned, the oblique projection amplifies numerical errors, so we solve Eq.~\eqref{eq:blockpm_smallproblem} with the pseudoinverse regularization described in Appendix~\ref{app:blockpm}.

The Ritz coefficient matrices are then used to de-mix the propagated blocks, assigning each member to one approximate eigenpair. The resulting states are truncated back to $\chi$ and normalized before the next iteration. We stop when the leading Ritz values remain stable within the prescribed tolerance.

The computational cost is dominated by the tensor-network operations rather than the $k\times k$ eigenproblem, and most of it sits in the de-mixing sums; Appendix~\ref{app:blockpm:cost} quantifies the breakdown and the measures that reduce it.

When only the eigenvalues are needed, we do not assign individual block members to separate eigenvectors at all: the blocks are kept in an orthonormal basis of the leading invariant subspace, which is better conditioned near a degeneracy and leaves the Ritz values unchanged (Appendix~\ref{app:blockpm:schur}). We use this mode for the late-time spectral measurements in Section~\ref{sec:results}.

Finally, we validate the implementation wherever $\mathcal{E}$ is small enough to be constructed explicitly and diagonalized exactly: the block method reproduces the leading dense spectra to between $10^{-8}$ and $10^{-13}$ (Appendix~\ref{app:blockpm:validation}).

\subsection{Selecting and tracking the physical branch}
\label{sec:reach}

The block returns $k$ Ritz values at each evolution time, and the physical $\mu_0$ has to be identified among them. We take the eigenvalue of largest modulus, and use continuity with the preceding rung to break the tie when two moduli agree to within one per cent. Such ties are produced by the conformal spectrum itself: as $T$ grows the moduli of the tower states approach a common value, which is the emergent dual unitarity of Section~\ref{sec:cft:temporal}. The subleading eigenvalues are then ordered by their phase distance from $\mu_0$, the order in which the conformal tower arranges them.

Along a ladder of evolution times the converged block at one time seeds the next, so consecutive rungs start from the same leading subspace and the selection stays on one branch as the spectrum tightens. When the leading block fills with tower members and part of the operator content falls outside it, the block size is escalated in steps ($k=4\to6\to8$), warm-seeding each larger solve with the converged smaller one, so the extra cost is paid only where the spectrum demands it. The safeguards protecting a ladder built on these rules, and the conditioning diagnostics monitored alongside it, are collected in Appendices~\ref{app:selector} and~\ref{app:conditioning}.
